\begin{problem}[Refinement invariance]\label{prob:vi20-p11}
Show refining an affine cover and using restricted transition maps reconstructs the same scheme.
\end{problem}

\ifdefined\FullProblemDossiers
\paragraph{Why this problem matters.}
Atlases are not unique, so reconstruction must be stable under refinement.

\paragraph{Useful ideas before starting.}
Both gluings map chartwise to the original scheme.

\paragraph{Guided hints.}
\begin{enumerate}
\item Both gluings map chartwise to the original scheme.
\end{enumerate}

% VI20_FULL_SOLUTION_REFINEMENT P11
\begin{solution}
Let $\{U_i\}$ be an affine cover of a scheme $X$, and let $\{V_\alpha\}$ be a refinement, so each $V_\alpha$ is contained in some $U_{i(\alpha)}$.  Give the refined cover the transition maps obtained by restricting the identity transitions on the original intersections.

Gluing the original cover reconstructs $X$: the chart inclusions $U_i\hookrightarrow X$ are compatible and induce an isomorphism from the original gluing to $X$.  The same argument applies to the refined cover: the inclusions $V_\alpha\hookrightarrow X$ are compatible and induce an isomorphism from the refined gluing to $X$.  Composing one reconstruction isomorphism with the inverse of the other gives an isomorphism between the two gluings.

This isomorphism is canonical because it is uniquely determined on each refined chart by the inclusion into the containing original chart.  Thus refinement changes the atlas but not the reconstructed scheme.  In particular, adding smaller affine charts or replacing a chart by an affine cover does not alter the global object.
\end{solution}


\paragraph{What happened mathematically.}
Changing coordinates or adding smaller charts does not change the underlying scheme.

\paragraph{Extensions.}
Use common refinements to compare two atlases.

\paragraph{Provenance.}
\textbf{New canonical synthesis; no numbered legacy Problem is claimed.}
\else
\paragraph{Provenance.} \textbf{New canonical synthesis; no numbered legacy Problem is claimed.}
\fi
